\documentclass{article}
\usepackage[solutions]{../../mystyle/Style}
\usepackage[margin=1in]{geometry}
\usepackage[colorlinks=true, linkcolor=blue, urlcolor=blue]{hyperref}

\title{LaTeX Reference and Snippet Examples}
\author{\defaultauthor}
\date{\today}

\begin{document}

\maketitle
\tableofcontents
\newpage

% =================================================
\section{Document Structure}
% =================================================

This is a section.

Text can contain inline mathematics such as
$x^2+y^2=z^2$ or \(e^{i\pi}+1=0\).

\subsection{A Subsection}

This is a subsection.

\subsubsection{A Subsubsection}

This is a smaller heading.


% =================================================
\section{Theorem Environments}
% =================================================

\begin{definition}
    A group is a set equipped with an associative
    binary operation, an identity element, and inverses.
\end{definition}

\begin{theorem}
    Every subgroup of a cyclic group is cyclic.
\end{theorem}

\begin{proof}
    Let $G=\langle g\rangle$ and let $H\leq G$.

    If $H=\{e\}$, it is cyclic.

    Otherwise, let $n$ be the smallest positive integer
    such that $g^n\in H$.

    By division with remainder, every $g^k\in H$
    is a power of $g^n$.

    Hence $H=\langle g^n\rangle$.
\end{proof}

\begin{lemma}
    An identity element of a group is unique.
\end{lemma}

\begin{proposition}
    Every cyclic group is abelian.
\end{proposition}

\begin{corollary}
    Every subgroup of a cyclic group is abelian.
\end{corollary}

\begin{example}
    The additive group $(\Z,+)$ is cyclic,
    generated by $1$.
\end{example}

\begin{remark}
    These statements use the theorem formatting
    supplied by \texttt{Style.sty}.
\end{remark}


% =================================================
\section{Math Modes and Equations}
% =================================================

\subsection{Inline Math}

Inline mathematics:

\[
    a^2+b^2=c^2
\]

You can also write $a^2+b^2=c^2$ directly
within a sentence.

\subsection{Display Math}

\[
    \frac{a+b}{c}=\sqrt{x^2+1}
\]

\subsection{Numbered Equations}

\begin{equation}
    \sum_{i=1}^{n}i=\frac{n(n+1)}{2}
    \label{eq:sample}
\end{equation}

\subsection{Unnumbered Equations}

\begin{equation*}
    E=mc^2
\end{equation*}

\subsection{Aligned Equations}

\begin{align*}
    (x+1)^2 &= x^2+2x+1 \\
            &= x(x+2)+1
\end{align*}

\subsection{Aligned Within Display Math}

\[
\begin{aligned}
    x^2 &= 4 \\
    x &= \pm 2
\end{aligned}
\]

\subsection{Text Within Math}

\[
    f(x)=x^2
    \quad \text{for all } x\in\R
\]

\subsection{Spacing}

\[
    a\,b
    \qquad
    a\:b
    \qquad
    a\;b
\]

\[
    x=2 \quad \text{and} \quad y=3
\]

\[
    x=2 \qquad y=3
\]


% =================================================
\section{Sets and Mathematical Commands}
% =================================================

\subsection{Number Sets}

\[
    \N\subseteq\Z\subseteq\Q\subseteq\R\subseteq\C
\]

\[
    \F
\]

\subsection{Sets}

\[
    A=\set{1,2,3}
\]

\[
    B=\set{x\in\R\mid x>0}
\]

\subsection{Absolute Values}

\[
    \abs{-5}=5
\]

\[
    \abs*{\frac{x}{y}}
\]

\subsection{Norms}

\[
    \norm{x}
\]

\[
    \norm*{\frac{x}{y}}
\]

\subsection{Ceiling and Floor}

\[
    \ceil{2.3}=3
\]

\[
    \floor{2.7}=2
\]

\subsection{Inner Products}

\[
    \inner{u}{v}
\]

\subsection{Vectors and Transposes}

\[
    \vect{v}
\]

\[
    A\transpose
\]

\subsection{Greek Letters}

\[
    \alpha,\quad
    \beta,\quad
    \gamma,\quad
    \delta,\quad
    \eps,\quad
    \theta,\quad
    \lambda,\quad
    \pi,\quad
    \sigma,\quad
    \omega
\]

\subsection{Mathematical Alphabets}

\[
    \mathbb{R},\quad
    \mathcal{F},\quad
    \mathfrak{g}
\]

\subsection{Common Operators}

\[
    \ker(f)
\]

\[
    \im(f)
\]

\[
    \rank(A)
\]

\[
    \Span\set{v_1,v_2}
\]

\[
    \Tr(A)
\]

\[
    \Var(X)
\]

\[
    \Cov(X,Y)
\]

\[
    \diag(A)
\]

\[
    \Res(f)
\]

\subsection{Definitions and Equality}

\[
    x\defeq 3
\]

\[
    3\eqdef x
\]

\subsection{Inverse Functions}

\[
    \phiinv
\]


% =================================================
\section{Calculus}
% =================================================

\subsection{Summations}

\[
    \sum_{i=1}^{n}i
\]

\[
    \sum_{n=1}^{\infty}\frac{1}{n^2}
\]

\subsection{Products}

\[
    \prod_{i=1}^{n}i=n!
\]

\subsection{Limits}

\[
    \lim_{n\to\infty}\frac{1}{n}=0
\]

\[
    \lim_{x\to 0}\frac{\sin x}{x}=1
\]

\subsection{Integrals}

\[
    \int_0^1 x^2\,dx=\frac13
\]

\[
    \int_a^b f(x)\,dx
\]


% =================================================
\section{Matrices and Delimiters}
% =================================================

\subsection{Parentheses}

\[
    \left(\frac{a}{b}\right)
\]

\subsection{Brackets}

\[
    \left[\frac{a}{b}\right]
\]

\subsection{Parenthesis Matrix}

\[
    \begin{pmatrix}
        a & b \\
        c & d
    \end{pmatrix}
\]

\subsection{Bracket Matrix}

\[
    \begin{bmatrix}
        1 & 2 \\
        3 & 4
    \end{bmatrix}
\]

\subsection{Piecewise Functions}

\[
    f(x)=
    \begin{cases}
        x^2 & \text{if }x\geq 0 \\
        -x & \text{if }x<0
    \end{cases}
\]


% =================================================
\section{Lists}
% =================================================

\subsection{Enumerate}

\begin{enumerate}
    \item First item
    \item Second item
    \item Third item
\end{enumerate}

\subsection{Itemize}

\begin{itemize}
    \item First bullet
    \item Second bullet
    \item Third bullet
\end{itemize}


% =================================================
\section{Homework Environments}
% =================================================

\subsection{Homework Information}

Your style defines the following commands:

\begin{verbatim}
\sethwinfo
    {\defaultauthor}
    {Math 113}
    {Homework 8}
    {November 20, 2026}

\makehwtitle
\end{verbatim}

These generate your homework heading.

\subsection{Problems}

\begin{problem}[Group Theory][10]
    Show that $\Z$ is an abelian group under addition.
\end{problem}

\subsection{Solutions}

\begin{solution}
    Addition in $\Z$ is associative and commutative.

    The identity is $0$.

    For every $n\in\Z$, the inverse is $-n$.
\end{solution}

\subsection{Subproblems}

\begin{subproblem}
    \item Check associativity.
    \item Check the identity.
    \item Check inverses.
\end{subproblem}


% =================================================
\section{Utility Commands}
% =================================================

\subsection{TODO}

\TODO{Review this section later.}

\subsection{Highlighting}

This is \highlight{important information}.

You can also use:

\highlight[green!25]{Green highlighting}

\subsection{Cross-References}

Equation~\ref{eq:sample} was defined earlier.

\subsection{Input Files}

You can combine other documents using:

\begin{verbatim}
\input{Lecture01}
\input{Lecture02}
\input{Lecture03}
\end{verbatim}

Each file's contents are inserted into the
current document.

\subsection{Table of Contents}

The command:

\begin{verbatim}
\tableofcontents
\end{verbatim}

automatically creates a table of contents
from your section headings.


% =================================================
\section{VS Code Snippet Reference}
% =================================================

\subsection{Document Templates}

\begin{itemize}
    \item \texttt{latexdoc} --- Basic standalone LaTeX document
    \item \texttt{mathdoc} --- Math document using \texttt{Style.sty}
    \item \texttt{basiclecture} --- Standalone lecture with title, author, and date
    \item \texttt{masterdoc} --- Master lecture document with a table of contents
    \item \texttt{lecturedoc} --- Individual lecture compatible with \texttt{masterdoc} through \texttt{subfiles}
    \item \texttt{hwdoc} --- Standalone homework document with problems and solutions
\end{itemize}

\subsection{Theorem Shortcuts}

\begin{itemize}
    \item \texttt{thm} --- Theorem
    \item \texttt{def} --- Definition
    \item \texttt{lem} --- Lemma
    \item \texttt{prop} --- Proposition
    \item \texttt{cor} --- Corollary
    \item \texttt{ex} --- Example
    \item \texttt{rem} --- Remark
    \item \texttt{prf} --- Proof
\end{itemize}

\subsection{Homework Shortcuts}

\begin{itemize}
    \item \texttt{prob} --- Problem
    \item \texttt{subprob} --- Subproblem
    \item \texttt{sol} --- Solution
\end{itemize}

\subsection{Equation Shortcuts}

\begin{itemize}
    \item \texttt{inline} --- Inline math
    \item \texttt{display} --- Display math
    \item \texttt{eq} --- Equation
    \item \texttt{align} --- Aligned equations
    \item \texttt{frac} --- Fraction
    \item \texttt{sqrt} --- Square root
\end{itemize}

\subsection{Math Shortcuts}

\begin{itemize}
    \item \texttt{sum} --- Summation
    \item \texttt{prod} --- Product
    \item \texttt{lim} --- Limit
    \item \texttt{int} --- Integral
    \item \texttt{set} --- Set
    \item \texttt{setb} --- Set builder
    \item \texttt{abs} --- Absolute value
    \item \texttt{norm} --- Norm
    \item \texttt{ceil} --- Ceiling
    \item \texttt{floor} --- Floor
    \item \texttt{inner} --- Inner product
\end{itemize}

\subsection{Formatting Shortcuts}

\begin{itemize}
    \item \texttt{bf;} --- Bold text
    \item \texttt{it;} --- Italic text
    \item \texttt{em;} --- Emphasized text
    \item \texttt{ul;} --- Underlined text
    \item \texttt{col;} --- Colored text
    \item \texttt{mbf;} --- Bold mathematical letters
    \item \texttt{bsy;} --- Bold mathematical symbols
    \item \texttt{mat} --- Parenthesis matrix
    \item \texttt{bmat} --- Bracket matrix
    \item \texttt{cases} --- Piecewise function
    \item \texttt{paren} --- Parentheses
    \item \texttt{bracket} --- Brackets
    \item \texttt{text} --- Text in math
    \item \texttt{bb} --- Blackboard bold
    \item \texttt{cal} --- Calligraphic
    \item \texttt{frak} --- Fraktur
\end{itemize}

\subsection{Document Organization Shortcuts}

\begin{itemize}
    \item \texttt{env} --- Generic environment
    \item \texttt{enum} --- Numbered list
    \item \texttt{item} --- Bulleted list
    \item \texttt{sec} --- Section
    \item \texttt{subsec} --- Subsection
    \item \texttt{input} --- Input another file
\end{itemize}

\end{document}